\begin{proof}[Proof of \cref{obj:prop:homogeneous-reduction}]
Recall the three covariance loads. For \(i,j\in O_n\) and an assignment-probability vector \(p=(p_k)_{k\in I_n}\) with \(0<p_k<1\),
\[
r_{ij}^1(G_n,p)
=\mathbf 1\{|S_{ij}(G_n)|>0\}
\left(\prod_{k\in S_{ij}(G_n)}p_k^{-1}-1\right),
\qquad
r_{ij}^0(G_n,p)
=\mathbf 1\{|S_{ij}(G_n)|>0\}
\left(\prod_{k\in S_{ij}(G_n)}(1-p_k)^{-1}-1\right),
\]
\[
r_{ij}^{10}(G_n)=\mathbf 1\{|S_{ij}(G_n)|>0\}.
\]
Throughout the proof, \(p\) denotes the probability vector and \(\rho\in(0,1)\) denotes the common homogeneous coordinate, so \(p_k=\rho\) for every \(k\in I_n\). We write \(\rho^{-\ell}\) as shorthand for \((\rho^{\ell})^{-1}\), and similarly for \((1-\rho)^{-\ell}\). We also adopt throughout the reciprocal convention
\[
1/0=0,
\]
so that the symbol \(1/n\), with \(n=\lvert O_n\rvert\), is defined for every finite outcome-unit set: it is the ordinary reciprocal when \(O_n\neq\varnothing\) and equals \(0\) when \(O_n=\varnothing\). Every identity displayed below is asserted under this convention, so the degenerate case \(O_n=\varnothing\) needs no separate hypothesis.

\begin{enumerate}
\item For every \(k\in I_n\), homogeneity and \(\rho\in(0,1)\) give
\[
0<p_k=\rho<1,
\]
so all exposure probabilities \(\pi_i^1(p)=\prod_{k\in N_i(G_n)}p_k\) and \(\pi_i^0(p)=\prod_{k\in N_i(G_n)}(1-p_k)\) are strictly positive and the reciprocals above are defined.
% lean: CausalSmith.Experimentation.BipartiteMinimaxDesign.homogeneous_reduction

\item Fix \(i,j\in O_n\). If \(\lvert S_{ij}(G_n)\rvert>0\), substituting \(p_k=\rho\) into the displayed definition of the treated load yields
\[
r_{ij}^1(G_n,p)
=\prod_{k\in S_{ij}(G_n)}\rho^{-1}-1
=(\rho^{\lvert S_{ij}(G_n)\rvert})^{-1}-1
=\rho^{-\lvert S_{ij}(G_n)\rvert}-1,
\]
because the product has \(\lvert S_{ij}(G_n)\rvert\) equal factors. If \(\lvert S_{ij}(G_n)\rvert=0\), the indicator in that definition vanishes and \(r_{ij}^1(G_n,p)=0\).
% lean: CausalSmith.Experimentation.BipartiteMinimaxDesign.homogeneous_reduction

\item The same split applies to the control load, with \(\rho\) replaced by \(1-\rho\). When
\(\lvert S_{ij}(G_n)\rvert>0\),
\[
r_{ij}^0(G_n,p)
=\prod_{k\in S_{ij}(G_n)}(1-\rho)^{-1}-1
=((1-\rho)^{\lvert S_{ij}(G_n)\rvert})^{-1}-1
=(1-\rho)^{-\lvert S_{ij}(G_n)\rvert}-1;
\]
when \(\lvert S_{ij}(G_n)\rvert=0\), it equals \(0\).
% lean: CausalSmith.Experimentation.BipartiteMinimaxDesign.homogeneous_reduction

\item It remains to prove the variance-scale identity. Write, for \(i\in O_n\),
\[
\delta_i^1=Y_i^1-\mu_1,\qquad \delta_i^0=Y_i^0-\mu_0,
\qquad
A_i=\frac{T_i(Z)}{\pi_i^1(p)}-1,
\qquad
B_i=\frac{C_i(Z)}{\pi_i^0(p)}-1,
\]
for the centered potential-outcome deviations and the centered inverse-probability weights, so that \cref{obj:def:first-order-linearization} reads \(\eta_i(p,Z)=A_i\delta_i^1-B_i\delta_i^0\). Under \cref{obj:ass:independent-heterogeneous-bernoulli} the coordinates of \(Z\) are independent, so
\[
\mathbb E_p\!\left[T_i(Z)\right]=\prod_{k\in N_i(G_n)}p_k=\pi_i^1(p),
\qquad
\mathbb E_p\!\left[C_i(Z)\right]=\prod_{k\in N_i(G_n)}(1-p_k)=\pi_i^0(p),
\]
both strictly positive by the coordinate bounds above; hence \(\mathbb E_p[A_i]=\mathbb E_p[B_i]=0\) and
\[
\mathbb E_p\!\left[\eta_i(p,Z)\right]=0 .
\]
% lean: CausalSmith.Experimentation.BipartiteMinimaxDesign.linScore_mean_zero

Unfolding \cref{obj:def:variance-scale} and applying the finite-design variance formula to the linear combination \(\sum_{i\in O_n}\frac1n\eta_i(p,Z)\) gives
\[
\sigma_{G_n,p}^2(Y)
= n\sum_{i\in O_n}\sum_{j\in O_n}
 \frac1n\cdot\frac1n\,
 \operatorname{Cov}_p\!\left(\eta_i(p,Z),\eta_j(p,Z)\right).
\]
The centering just proved turns each covariance into the corresponding raw pair moment. If \(O_n\neq\varnothing\), then \(n>0\) and the scalar identity \(n\cdot(1/n)^2=1/n\) applies; if \(O_n=\varnothing\), both the displayed right-hand side and the double sum below are empty, hence \(0\), and the two sides again agree under the reciprocal convention. In either case
\[
\sigma_{G_n,p}^2(Y)
=\frac1n\sum_{i\in O_n}\sum_{j\in O_n}
\mathbb E_p\!\left[\eta_i(p,Z)\eta_j(p,Z)\right].
\]
% lean: CausalSmith.Experimentation.BipartiteMinimaxDesign.varScale_pair_moments

Next evaluate those pair moments. For any random variables \(U,V\) with \(\mathbb E_p[U]=\theta_U\neq0\) and \(\mathbb E_p[V]=\theta_V\neq0\),
\(\mathbb E_p[(U/\theta_U-1)(V/\theta_V-1)]=\mathbb E_p[UV]/(\theta_U\theta_V)-1\). Independence gives
\(\mathbb E_p[T_i(Z)T_j(Z)]=\prod_{k\in N_i(G_n)\cup N_j(G_n)}p_k\), while
\(\pi_i^1(p)\pi_j^1(p)=\bigl(\prod_{k\in N_i(G_n)\cup N_j(G_n)}p_k\bigr)\bigl(\prod_{k\in S_{ij}(G_n)}p_k\bigr)\), so
\[
\mathbb E_p\!\left[A_iA_j\right]
=\prod_{k\in S_{ij}(G_n)}p_k^{-1}-1
=r_{ij}^1(G_n,p).
\]
The same argument with \(1-p_k\) in place of \(p_k\) gives
\[
\mathbb E_p\!\left[B_iB_j\right]
=\prod_{k\in S_{ij}(G_n)}(1-p_k)^{-1}-1
=r_{ij}^0(G_n,p),
\]
with the empty-overlap case reducing to \(1-1=0\), the value selected by the indicators in \(r_{ij}^1\) and \(r_{ij}^0\). For the mixed moments, if \(S_{ij}(G_n)\neq\varnothing\), then \(T_i(Z)C_j(Z)=0\) identically, since a shared intervention cannot be both treated and untreated, so \(\mathbb E_p[A_iB_j]=-1\). If \(S_{ij}(G_n)=\varnothing\), then the two exposure events involve disjoint assignment coordinates, and independence gives \(\mathbb E_p[T_i(Z)C_j(Z)]=\pi_i^1(p)\pi_j^0(p)\), hence \(\mathbb E_p[A_iB_j]=0\). In both cases, and likewise with the roles of the two outcome units exchanged,
\[
\mathbb E_p\!\left[A_iB_j\right]=\mathbb E_p\!\left[B_iA_j\right]=-r_{ij}^{10}(G_n).
\]
Expanding \(\eta_i(p,Z)\eta_j(p,Z)=A_iA_j\delta_i^1\delta_j^1-A_iB_j\delta_i^1\delta_j^0-B_iA_j\delta_i^0\delta_j^1+B_iB_j\delta_i^0\delta_j^0\) and inserting the four displayed moments gives
\[
\mathbb E_p\!\left[\eta_i(p,Z)\eta_j(p,Z)\right]
=r_{ij}^1(G_n,p)\delta_i^1\delta_j^1
+r_{ij}^0(G_n,p)\delta_i^0\delta_j^0
+r_{ij}^{10}(G_n)\left(\delta_i^1\delta_j^0+\delta_i^0\delta_j^1\right).
\]
% lean: CausalSmith.Experimentation.BipartiteMinimaxDesign.linScore_pair_moment

Finally, \(S_{ij}(G_n)=S_{ji}(G_n)\), so \(r_{ij}^{10}(G_n)=r_{ji}^{10}(G_n)\); relabelling \((i,j)\mapsto(j,i)\) in the double sum gives
\[
\sum_{i\in O_n}\sum_{j\in O_n}r_{ij}^{10}(G_n)\delta_i^0\delta_j^1
=\sum_{i\in O_n}\sum_{j\in O_n}r_{ij}^{10}(G_n)\delta_i^1\delta_j^0 .
\]
The two mixed contributions therefore combine into a single doubled term. Substituting the pair-moment identity into the pair expansion of \(\sigma_{G_n,p}^2(Y)\) yields
\[
\sigma_{G_n,p}^2(Y)
=
\frac{1}{n}\sum_{i\in O_n}\sum_{j\in O_n}
\left[
r_{ij}^1(G_n,p)(Y_i^1-\mu_1)(Y_j^1-\mu_1)
+r_{ij}^0(G_n,p)(Y_i^0-\mu_0)(Y_j^0-\mu_0)
+2r_{ij}^{10}(G_n)(Y_i^1-\mu_1)(Y_j^0-\mu_0)
\right],
\]
as asserted.
% lean: CausalSmith.Experimentation.BipartiteMinimaxDesign.varScale_homogeneous_formula
\end{enumerate}
\end{proof}
